% !TEX root = thesis.tex
% =====================================================================
%  CHAPTER 7 — CONCLUSION
%  Zero-Shot Nutrient and Quantity Association using
%  OCR-Guided Semantic Graph Matching
%  Moustafa Hamada | THD + USB
%
%  Build: scrbook, XeLaTeX + BibTeX. Cross-refs: Chapter~\ref{} / Section~\ref{}.
%
%  This file DEFINES \label{chap:conclusion}, which resolves the forward
%  reference made from the Introduction (Chapter 1). Placed last in the
%  document, after Chapters 1--6.
%
%  BACK REFERENCES (labels already defined in the other chapter files):
%    chap:introduction (Ch1), chap:experiments (Ch6)
%
%  This chapter is high-level and adds NO new citations (it refers back to the
%  chapters where each component and result was introduced).
%
%  NOTE: \mcode{} is defined inside Chapter 4. In a full front-to-back build it
%  would be available here, but this chapter is kept free of \mcode so it also
%  compiles standalone; edge types and modes are named in plain words instead.
%
%  All figures for the four-field results (0.4304 / 0.4945 / 0.3597, the OCR
%  jump, the runtime split) are stated verbatim from the locked runs reported
%  in Chapter~\ref{chap:experiments}.
%
%  LABELS IN THIS FILE:
%    chap:conclusion
%    sec:concl-findings, sec:concl-limitations, sec:concl-future,
%    sec:concl-closing
% =====================================================================

\chapter{Conclusion}\label{chap:conclusion}

This thesis set out to reconstruct structured nutritional records from images of
supplement and food labels, under two demands that usually pull against each
other. The method had to work zero-shot, on labels, layouts, and languages it had
never seen, and it had to stay open to inspection at every stage rather than hand
the task to a single opaque model. The answer proposed and tested here is a
modular pipeline that reads the label with OCR, classifies each token into a
semantic role without training on the target labels, and then binds the tokens
into four-field tuples
$\langle\text{nutrient},\,\text{quantity},\,\text{unit},\,\text{context}\rangle$
by matching an explicit graph whose edges follow closed-form geometric and
structural rules. This chapter reviews what the pipeline achieved against the
questions set in Chapter~\ref{chap:introduction}, states where it falls short, and
points to the work that would take it further.

\section{Summary of Findings}\label{sec:concl-findings}

The experiments in Chapter~\ref{chap:experiments} answered each of the five
questions by ablation, changing one stage at a time. The reading stage proved to
be the single largest lever. Replacing EasyOCR with PaddleOCR PP-OCRv5, with every
later stage held fixed, lifted the four-field tuple F1 from 0.1773 to 0.4231, and
the gain came almost entirely from reading the numbers and units correctly rather
than from finding more nutrient names. Because OCR sets the ceiling for everything
after it, PaddleOCR was fixed for the rest of the work.

The association stage was the next-largest factor, and it is where the
contribution sits. The enriched, geometry-aware graph nearly doubled the
four-field F1 of the naive geometric graph for both learned classifiers, and the
gain was structural rather than a matter of better reading: it appeared almost
entirely on the multi-context labels, the ones that report the same nutrient under
more than one reference basis, where the header-scope relation keeps the value
columns apart and the simpler graph cannot. Inside the enriched graph, geometric
row edges were clearly the right mechanism. The geometric overlap mode was best,
the rank-based modes were far weaker --- a single missed or fused token shifts the
ordinal count and breaks the alignment --- and merging rank edges into the
geometric ones only added wrong links. This confirmed the central design choice:
on labels distorted by OCR, physical alignment is a more reliable basis for the
row relation than ordinal rank.

The classifier was the smallest of the three factors, because it acts only through
nutrient detection while the shared cascade reads quantity, unit, and context the
same way in every run. Among the zero-shot options, GLiNER-biomed gave the best
interpretable tuple score, and it did so without any hand-built lexicon, which
matters for a method that claims to generalise beyond a fixed vocabulary.

Taken together, these choices give the central result of the thesis. The best
fully interpretable pipeline --- PaddleOCR, GLiNER, and the enriched graph with
geometric overlap edges --- reaches a four-field tuple F1 of 0.4304, with every
stage open to inspection. Confining a vision--language model to the association
step, so that it binds the same token table the graph receives, raises the score
to 0.4945; with the model doing the binding, the embedding classifier overtakes
GLiNER, since its higher recall becomes an advantage rather than a source of false
tuples. Both of these beat a vision--language model run end to end, which reaches
only 0.3597. This is the answer the central question asked for: structuring the
problem before any field is bound is worth more than handing the whole task to a
single model, and the extra points the confined model adds are paid for in
interpretability and in time --- the interpretable pipeline runs in about four and
a half seconds per label, the hybrid in about fourteen. Where the task needs a
result that can be traced and checked, the interpretable pipeline is the answer;
where the highest score matters more than the audit trail, the confined model is.

\section{Limitations}\label{sec:concl-limitations}

Two limits bound every configuration, and both arise before the binding step, so
no choice of graph or model repairs them. The first is a recall ceiling. About a
third of the gold tuples --- 314 of the 866 --- are lost because their nutrient is
never read, whether the OCR corrupts the text or the classifier rejects a genuine
nutrient. Once a nutrient is missing, none of its values can be recovered, and
this loss passes through every run unchanged. The pipeline can only bind the
tokens it is handed.

The second is the quantity field, which is the weakest of the four. After a
nutrient is matched, its number is wrong more often than its unit or its context,
and the errors concentrate on dense tables that place two value columns side by
side, where the enrichment stage can assign a value to the wrong column and a
per-serving number is bound where the per-100-gram number belongs. This is a fault
in how the columns are separated, not in the binding itself, so it is best fixed
upstream of the graph.

Beyond these, the association step is deliberately simple: it works greedily from
each nutrient and does not enforce a single global assignment, so it can leave
some conflicts between competing links unresolved. Labels that are not laid out as
tables, such as those printed as running paragraphs, fall outside what the graph
is built to read and are only partly recovered by a regular-expression fallback.
And the evaluation rests on a corpus of 57 labels; although the labels are
multilingual and varied, a larger and more diverse set --- more languages, more
scripts, and more unusual layouts --- would test the zero-shot claim more severely
than the present corpus can.

\section{Future Work}\label{sec:concl-future}

The clearest gains would come from lifting the recall ceiling, since the reading
stage is both the largest lever and the source of the largest loss. Stronger or
better-tuned OCR, wider correction rules for corrupted nutrient names, or a second
recognition pass aimed at the tokens the current engine drops would each recover
tuples that the rest of the pipeline is ready to bind but never receives.

The quantity weakness points to a specific fix upstream of the graph. Because the
errors come from values being placed in the wrong column on dense multi-column
tables, a more robust column-separation step in the enrichment stage --- one that
is less easily disturbed by a missing or extra token --- would attack the weakest
field at its source, without changing the association logic.

On the binding step itself, replacing the greedy local search with a global
one-to-one assignment would resolve the conflicts the current method leaves open,
while keeping the graph interpretable. A different line would keep the accuracy of
the confined vision--language model but recover some of the transparency it costs:
rather than letting the model bind the tokens freely, it could be used only to
re-rank or confirm the candidate links the graph already proposes, so that every
accepted link still traces back to an explicit edge. This would sit between the
interpretable and the hybrid pipeline, trading a little of the model's score for
an audit trail.

Finally, a broader evaluation would strengthen the zero-shot argument. Testing the
pipeline on a larger corpus with more languages, more scripts, and more layout
styles would show how far the closed-form, training-free design carries beyond the
labels seen here, and would expose any layout family the current rules do not yet
handle.

\section{Concluding Remarks}\label{sec:concl-closing}

The wider point of this thesis is that a structured, interpretable pipeline is a
viable way to read nutrition labels, not a compromise forced by the absence of a
large model. The task was framed in Chapter~\ref{chap:introduction} as an empty
corner of the design space --- a method that is at once zero-shot, interpretable,
relational, and able to handle multilingual labels, which no existing approach
occupies. The pipeline built here fills that corner. Its association step
reconstructs four-field records from arbitrary, previously unseen labels without
any task-specific training, and every record it produces can be traced back
through an explicit graph to the tokens and rules that formed it. That a fully
interpretable pipeline beats an end-to-end model on the same task, and comes
within a few points of a version that gives up its transparency, suggests that
closed-form geometric reasoning over classified OCR tokens is worth taking
seriously as an alternative to opaque extraction --- one whose main cost is not
accuracy but the effort of building the structure that makes it legible.